% ============================================================
%  Preambolo aggiuntivo per la dispensa di Tecnologie Web
%  Estende unina_doc_class.cls (stile ripreso dalle Dispense APA)
% ============================================================

% ─── Font aggiuntivi ─────────────────────────────────────────
% \Mononoki e' richiamato dalla classe nelle macro \dir e \file
\newfontfamily\Mononoki{mononoki}[
  Path = fonts/, Extension = .ttf,
  UprightFont = *-Regular, BoldFont = *-Bold,
  ItalicFont  = *-Italic,  BoldItalicFont = *-BoldItalic]

% \Lora e' usata per simulare il rendering di default del browser
\newfontfamily\Lora{Lora}[
  Path = fonts/, Extension = .ttf,
  UprightFont = *-Regular, BoldFont = *-Bold,
  ItalicFont  = *-Italic,  BoldItalicFont = *-BoldItalic]

% FiraCode SENZA legature di programmazione.
% La classe dichiara \FiraCode con Contextuals=Alternate: in un testo
% che parla di HTML questo e' dannoso, perche' "<!--" verrebbe reso
% come "<!—" e "-->" come "—>", cioe' esattamente non quello che lo
% studente deve digitare. Qui si ridefinisce \FiraCode senza legature.
\newfontfamily\FiraCodeNL{FiraCode}[
  Path = fonts/, Extension = .ttf,
  UprightFont = *-Regular, BoldFont = *-Bold,
  Ligatures = NoCommon, Contextuals = NoAlternate]
\let\FiraCode\FiraCodeNL

% Il basicstyle dei listati va riallineato al font senza legature.
% Si disattiva anche mathescape (attivo nella classe): in un listato CSS
% un "$" e' un carattere legittimo (es. il selettore [href$='.it/']) e non
% deve aprire la modalita' matematica.
\lstset{basicstyle=\linespread{1.15}\footnotesize\FiraCodeNL, mathescape=false}

% ─── Etichette richieste dalla classe per l'elenco dei listati ─
\gdef\documentlsttitle{Elenco dei listati}
\gdef\documentlstlabel{Listato}

% ─── Colori dedicati all'evidenziazione del codice ───────────
\definecolor{codetag}{HTML}{1B4F9C}      % nomi dei tag
\definecolor{codeattr}{HTML}{9B2D8F}     % nomi degli attributi
\definecolor{codestring}{HTML}{1F7A5A}   % valori / stringhe
\definecolor{codecomment}{HTML}{8A8A8A}  % commenti
\definecolor{codekey}{HTML}{C2185B}      % parole chiave
\definecolor{codenum}{HTML}{B36B00}      % numeri e unita'
\definecolor{codehead}{HTML}{7A3E00}     % header HTTP

% ============================================================
%  LINGUAGGI PER LISTINGS
% ============================================================

% Le chiavi tag / tagstyle / markfirstintag / usekeywordsintag fanno
% parte dell'"aspetto" html di listings, che va caricato esplicitamente.
\lstloadaspects{html}

% ─── HTML ────────────────────────────────────────────────────
% Le stringhe sono delimitate SOLO dai doppi apici: in questo modo
% gli apostrofi presenti nel testo in italiano non aprono una stringa.
% Il commento va dichiarato DOPO il tag: se lo si dichiara prima, la
% sequenza "<!--" viene letta come apertura di un tag e la prima parola del
% commento prende lo stile del nome del tag.
\lstdefinelanguage{HTML5}{
  sensitive=false,
  morestring=[b]",
  tag=**[s]<>,
  usekeywordsintag=false,
  markfirstintag=true,
  morecomment=[s]{<!--}{-->},
}

% ─── CSS (usato nella seconda parte della dispensa) ──────────
\lstdefinelanguage{CSS3}{
  sensitive=false,
  alsoletter={-},
  morecomment=[s]{/*}{*/},
  morestring=[b]",
  morestring=[b]',
  morekeywords={
    color,background,background-color,background-image,font,font-size,
    font-family,font-weight,font-style,text-align,text-decoration,
    margin,margin-top,margin-bottom,margin-left,margin-right,
    padding,padding-top,padding-bottom,padding-left,padding-right,
    border,border-radius,border-color,border-width,border-style,
    width,height,max-width,min-width,max-height,min-height,
    display,position,top,left,right,bottom,float,clear,overflow,
    flex,flex-direction,justify-content,align-items,gap,grid,
    grid-template-columns,grid-template-rows,z-index,opacity,
    box-sizing,line-height,letter-spacing,content,cursor,transition},
}

% ─── JavaScript (usato nella seconda parte della dispensa) ───
\lstdefinelanguage{JS}{
  sensitive=true,
  morecomment=[l]{//},
  morecomment=[s]{/*}{*/},
  morestring=[b]",
  morestring=[b]',
  morestring=[b]`,
  morekeywords={
    var,let,const,function,return,if,else,for,while,do,switch,case,
    break,continue,new,delete,typeof,instanceof,this,class,extends,
    super,import,export,default,try,catch,finally,throw,async,await,
    yield,of,in,null,undefined,true,false},
}

% ─── JSON (package.json, payload delle API REST) ─────────────
\lstdefinelanguage{JSON}{
  sensitive=true,
  morestring=[b]",
  morecomment=[l]{//},
  literate=
    *{\{}{{{\color{codeattr}\{}}}{1}
     {\}}{{{\color{codeattr}\}}}}{1}
     {[}{{{\color{codeattr}[}}}{1}
     {]}{{{\color{codeattr}]}}}{1}
     {:}{{{\color{codeattr}:}}}{1}
     {,}{{{\color{codeattr},}}}{1},
  morekeywords={true,false,null},
}

% ─── TypeScript (JavaScript + tipi) ──────────────────────────
\lstdefinelanguage{TS}[]{JS}{
  morekeywords={type,interface,enum,implements,readonly,public,private,
                protected,abstract,declare,namespace,as,is,keyof,infer,
                string,number,boolean,any,unknown,never,void,object,symbol},
}

% ─── Pug (motore di template) ────────────────────────────────
\lstdefinelanguage{Pug}{
  sensitive=true,
  morecomment=[l]{//-},
  morecomment=[l]{//},
  morestring=[b]",
  morestring=[b]',
  alsoletter={-},
  morekeywords={doctype,html,head,body,title,script,link,meta,
                block,extends,include,each,in,if,else,unless,
                h1,h2,h3,p,ul,li,ol,a,footer,header,nav,main,
                article,section,aside,div,span,form,label,input,
                button,table,tr,td,th,img},
}

% ─── Trascrizioni di terminale ───────────────────────────────
\lstdefinelanguage{Shell}{
  sensitive=true,
  morecomment=[l]{\#},
  morestring=[b]",
  morestring=[b]',
  morekeywords={npm,node,npx,git,cd,ls,mkdir,curl,docker,nodemon,
                install, init, start, run, test, build},
}

% ─── HTTP (messaggi grezzi di richiesta/risposta) ────────────
\lstdefinelanguage{HTTP}{
  sensitive=true,
  alsoletter={/.-},
  morestring=[b]",
  keywords=[1]{GET,POST,PUT,DELETE,PATCH,HEAD,OPTIONS,TRACE,CONNECT},
  keywords=[2]{HTTP/1.0,HTTP/1.1,HTTP/2,HTTP/3},
  keywords=[3]{Host,User-Agent,Accept,Accept-Language,Accept-Encoding,
    Connection,Content-Type,Content-Length,Content-Encoding,Cookie,
    Set-Cookie,Location,Server,Date,Cache-Control,Authorization,
    Referer,Last-Modified,ETag,If-None-Match},
}

% ============================================================
%  STILI DI LISTATO
% ============================================================

% literate={} annulla le sostituzioni dello stile FiraCodeStyle caricato dalla
% classe: quello stile trasforma "</", "/>", "<!--" e "-->" in simboli unici,
% e cosi' facendo impedisce a listings di riconoscere i tag di CHIUSURA (che
% restavano neri) e i commenti. Le legature sono comunque gia' disattivate nel
% font, quindi a video non cambia nulla se non la colorazione, ora corretta.
\lstdefinestyle{html}{
  language=HTML5,
  literate={},
  keywordstyle=\bfseries\color{codetag},   % nome del tag
  tagstyle=\color{codeattr},               % parentesi angolari e attributi
  stringstyle=\color{codestring},          % valori degli attributi
  commentstyle=\itshape\color{codecomment},
}

\lstdefinestyle{css}{
  language=CSS3,
  keywordstyle=\color{codetag},
  stringstyle=\color{codestring},
  commentstyle=\itshape\color{codecomment},
}

\lstdefinestyle{js}{
  language=JS,
  keywordstyle=\color{codekey},
  stringstyle=\color{codestring},
  commentstyle=\itshape\color{codecomment},
}

\lstdefinestyle{http}{
  language=HTTP,
  keywordstyle=[1]\bfseries\color{codekey},
  keywordstyle=[2]\bfseries\color{codetag},
  keywordstyle=[3]\color{codehead},
  stringstyle=\color{codestring},
}

\lstdefinestyle{plain}{
  language={},
  keywordstyle=\color{codekey},
  stringstyle=\color{codestring},
  commentstyle=\itshape\color{codecomment},
}

\lstdefinestyle{php}{
  language=PHP,
  keywordstyle=\color{codekey},
  stringstyle=\color{codestring},
  commentstyle=\itshape\color{codecomment},
  identifierstyle=\color{black},
}

\lstdefinestyle{json}{
  language=JSON,
  keywordstyle=\color{codekey},
  stringstyle=\color{codestring},
  commentstyle=\itshape\color{codecomment},
}

\lstdefinestyle{shell}{
  language=Shell,
  keywordstyle=\bfseries\color{codetag},
  stringstyle=\color{codestring},
  commentstyle=\itshape\color{codecomment},
}

\lstdefinestyle{ts}{
  language=TS,
  keywordstyle=\color{codekey},
  stringstyle=\color{codestring},
  commentstyle=\itshape\color{codecomment},
}

\lstdefinestyle{pug}{
  language=Pug,
  keywordstyle=\bfseries\color{codetag},
  stringstyle=\color{codestring},
  commentstyle=\itshape\color{codecomment},
}

\lstdefinestyle{sql}{
  language=SQL,
  keywordstyle=\bfseries\color{codekey},
  stringstyle=\color{codestring},
  commentstyle=\itshape\color{codecomment},
}

% Il titolo del listato (il nome del file) deve comparire solo se e' stato
% davvero specificato con la chiave "name". Non basta rendere vuoto il testo
% del titolo: listings emetterebbe comunque il box della didascalia, con il
% suo spazio verticale, lasciando un buco sotto i frammenti anonimi.
% Il blocco che stampa la didascalia e' protetto da \ifx\lst@caption\@empty,
% quindi si azzera \lst@caption quando manca il nome. (Nel documento non si
% usa mai la chiave "caption": quella "title" della classe la rende comunque
% inservibile.)
\makeatletter
\lstset{title={\color{primary}\fontsize{10}{12}\selectfont
               \sp@rtext{\lstname}{\ifilecode}}}
\lst@AddToHook{InitVars}{%
  \ifx\lst@intname\@empty
    \global\let\lst@caption\@empty
  \else
    \global\let\lst@caption\relax
  \fi}
\makeatother

% ============================================================
%  CORREZIONE DEI RIQUADRI DELLA CLASSE
% ============================================================
% Nella classe il titolo viene passato a tcolorbox come "title=#1":
% una virgola (o un "=") nel titolo verrebbe interpretata come
% separatore di chiavi. Racchiudendo il titolo in un gruppo il
% problema sparisce e si possono scrivere titoli in italiano normale.
\renewcommand{\info}[2]{\begin{infobox}{{#1}}#2\end{infobox}}
\renewcommand{\warning}[2]{\begin{warningbox}{{#1}}#2\end{warningbox}}
\renewcommand{\error}[2]{\begin{errorbox}{{#1}}#2\end{errorbox}}

% ============================================================
%  NUMERAZIONE DI ESEMPI ED ESERCIZI
% ============================================================
% Nella classe i contatori sono agganciati a "section" e l'etichetta
% e' "\thesection.\the<contatore>". Ma con secnumdepth=0 il contatore
% di sezione non viene mai incrementato, per cui si otterrebbe
% "Esempio 1.0.3". Li riaggancio quindi al capitolo.
\counterwithout{esempio}{section}
\counterwithin{esempio}{chapter}
\counterwithout{esercizio}{section}
\counterwithin{esercizio}{chapter}

% I \nopagebreak tengono il filetto e il titolo dell'esempio attaccati al
% contenuto che segue: senza, il titolo poteva restare da solo in fondo a una
% pagina, con il codice (che non si spezza) spostato alla pagina successiva.
% L'esempio viene composto in una scatola: se sta in una pagina la si mette
% intera (e passa alla pagina dopo se non entra), altrimenti la si "svuota"
% nella pagina e l'esempio può spezzarsi normalmente.
\newbox\wtesbox
\newcommand{\wtemetti}{%
  \ifdim\dimexpr\ht\wtesbox+\dp\wtesbox\relax>\textheight
    \unvbox\wtesbox
  \else
    \box\wtesbox
  \fi}
\renewenvironment{esempio}[1]{%
  \par\addvspace{0.8em}%
  \setbox\wtesbox\vbox\bgroup
  \refstepcounter{esempio}%
  \noindent\rule{\linewidth}{0.5pt}\par\nopagebreak%
  \vspace{0.15em}%
  \noindent{\color{primary!80!black}\textbf{Esempio \theesempio}}%
  \ifx\relax#1\relax\else%
    {\color{primary!80!black}\textbf{ --- #1}}%
  \fi\par\nopagebreak%
  \vspace{0.3em}%
}{%
  \par\vspace{0.3em}%
  \noindent\rule{\linewidth}{0.5pt}\par
  \egroup
  \wtemetti
  \par\vspace{0.8em}%
}

\renewenvironment{esercizio}[1]{%
  \par\addvspace{0.8em}%
  \setbox\wtesbox\vbox\bgroup
  \refstepcounter{esercizio}%
  \noindent\rule{\linewidth}{1.2pt}\vspace{1pt}\par%
  \noindent\rule{\linewidth}{0.4pt}\par%
  \vspace{0.15em}%
  \noindent{\color{accent!80!black}\textbf{Esercizio \theesercizio}}%
  \ifx\relax#1\relax\else%
    {\color{accent!80!black}\textbf{ --- #1}}%
  \fi\par%
  \vspace{0.3em}%
}{%
  \par\vspace{0.3em}%
  \noindent\rule{\linewidth}{0.4pt}\vspace{1pt}\par%
  \noindent\rule{\linewidth}{1.2pt}\par
  \egroup
  \wtemetti
  \par\vspace{0.8em}%
}

% ============================================================
%  MACRO INLINE
% ============================================================

% Codice inline breve (non spezzabile)
\def\code#1{{\FiraCode\color{primary!65!black}\mbox{#1}}}
% Codice inline lungo (spezzabile a fine riga)
\def\codel#1{{\FiraCode\color{primary!65!black}#1}}
% Tag HTML: \tg{p} produce <p>
\def\tg#1{\code{<#1>}}
% Attributo HTML
\def\attr#1{{\FiraCode\color{codeattr}\mbox{#1}}}
% Termine tecnico evidenziato (spezzabile, a differenza di \hi)
\def\term#1{{\color{accent!80!black}\LatoBlack #1}}
% URL mostrata come testo
\def\urlx#1{{\FiraCode\small\color{primary!70!black}\url{#1}}}

% ============================================================
%  FINESTRA "BROWSER" per mostrare il risultato del rendering
%  Uso: \begin{browser}{etichetta opzionale} ... \end{browser}
% ============================================================
\newcommand{\wtbrowserbefore}{0.35cm}
\newcommand{\wtbrowserafter}{0.6cm}
\newtcolorbox{browser}[1][]{%
  enhanced,
  colback=white,
  colframe=neutral!35!white,
  boxrule=0.8pt,
  arc=3pt,
  drop fuzzy shadow,
  before skip=\wtbrowserbefore, after skip=\wtbrowserafter,
  left=8pt, right=8pt, top=6pt, bottom=6pt,
  fonttitle=\footnotesize\Lato,
  colbacktitle=neutral!15!white,
  coltitle=neutral!35!black,
  title={%
    \tikz[baseline=-2.5pt]{%
      \fill[red-gradient-3] (0,0) circle (2.2pt);
      \fill[yellow-gradient-5] (8pt,0) circle (2.2pt);
      \fill[green-gradient-6] (16pt,0) circle (2.2pt);}%
    \hspace{8pt}\footnotesize #1},
  attach boxed title to top left={yshift*=-\tcboxedtitleheight},
  boxed title style={
    sharp corners=downhill, arc=3pt,
    colframe=neutral!35!white, boxrule=0.8pt,
  },
  % dentro la finestra si usa un font con grazie, come il default dei browser
  fontupper=\Lora\small,
}

% ============================================================
%  CONSOLE del browser, per mostrare l'output di uno script
%  Uso: \begin{console} riga1 \\ riga2 \end{console}
% ============================================================
\definecolor{consolebg}{HTML}{22252E}
\definecolor{consolefg}{HTML}{E4E6EE}
\definecolor{consolebar}{HTML}{31353F}

% Il contenuto e' VERBATIM: l'output di un programma contiene spesso
% caratteri speciali per LaTeX (_ \ ^ $ % &, ASCII art...) e gli a-capo
% devono essere preservati cosi' come sono.
% Come per lstlisting, questo ambiente non puo' stare dentro l'argomento
% di una macro: nei riquadri va usata la forma ambiente (warningbox...).
\newcommand{\wtconsolebefore}{0.35cm}
\newtcblisting{console}[1][Console]{%
  listing only,
  enhanced,
  colback=consolebg,
  colframe=consolebg,
  boxrule=0pt,
  arc=3pt,
  drop fuzzy shadow,
  before skip=\wtconsolebefore, after skip=0.6cm,
  left=8pt, right=8pt, top=6pt, bottom=6pt,
  fonttitle=\footnotesize\Lato,
  colbacktitle=consolebar,
  coltitle=consolefg!75!consolebg,
  title={\iterminal\hspace{6pt}\footnotesize #1},
  attach boxed title to top left={yshift*=-\tcboxedtitleheight},
  boxed title style={sharp corners=downhill, arc=3pt, boxrule=0pt,
                     colframe=consolebar},
  listing options={
    language={},
    basicstyle=\linespread{1.1}\footnotesize\FiraCodeNL\color{consolefg},
    backgroundcolor=\color{consolebg},
    keywordstyle={}, commentstyle={}, stringstyle={},
    numbers=none, frame=none,
    xleftmargin=0pt, xrightmargin=0pt,
    breaklines=true, columns=fixed,
    mathescape=false, escapechar={},
  },
}

% ============================================================
%  STILI TIKZ CONDIVISI DALLE FIGURE
% ============================================================
\tikzset{
  wtnode/.style={
    rounded corners=3pt, draw=primary!65!black, fill=primary!8!white,
    line width=0.8pt, inner sep=5pt, align=center, font=\Lato\small},
  wtnodeacc/.style={
    wtnode, draw=accent!60!black, fill=accent!8!white},
  wtnodewarn/.style={
    wtnode, draw=orange-gradient-6, fill=orange-gradient-1!45!white},
  wtnodeok/.style={
    wtnode, draw=green-gradient-6, fill=green-gradient-1!45!white},
  wtnodeplain/.style={
    rounded corners=2pt, draw=neutral!60!white, fill=neutral!8!white,
    line width=0.7pt, inner sep=4pt, align=center, font=\Lato\footnotesize},
  wtarrow/.style={
    -{Stealth[length=6pt,width=5pt]}, line width=0.9pt, draw=primary!70!black},
  wtarrowacc/.style={
    -{Stealth[length=6pt,width=5pt]}, line width=0.9pt, draw=accent!70!black},
  wtdash/.style={
    dashed, line width=0.7pt, draw=neutral!50!white},
  wtlab/.style={
    font=\Lato\scriptsize, text=neutral!15!black, inner sep=2.5pt,
    fill=white, align=center},
  wtbrace/.style={
    decorate, decoration={brace, amplitude=4pt, mirror}, draw=neutral!40!black,
    line width=0.7pt},
}

% Ambiente figura centrata con didascalia, usato in tutta la dispensa
\newcommand{\wtfig}[3][0.95]{%
  \begin{figure}[H]
    \centering
    \resizebox{#1\linewidth}{!}{#2}
    \caption{#3}
  \end{figure}}

% ============================================================
%  IL CODICE NON SI SPEZZA FRA DUE PAGINE
% ============================================================
% Ogni lstlisting viene racchiuso in una minipage. Una minipage e' una scatola
% indivisibile: se il listato non entra nello spazio rimasto in fondo alla
% pagina, passa per intero alla pagina successiva invece di essere tagliato.
% Il listato piu' lungo della dispensa (38 righe) sta comodamente in una pagina.
\AddToHook{env/lstlisting/before}{\wtattacca\noindent\begin{minipage}{\linewidth}}
\AddToHook{env/lstlisting/after}{\end{minipage}\par}

% ============================================================
%  CODICE E RISULTATO AFFIANCATI
% ============================================================
% Uso:
%   \begin{affianco}[0.55]      % frazione della riga occupata dal codice
%     \begin{lstlisting}...\end{lstlisting}   (uno o piu' listati)
%   \risultato
%     \begin{browser}[file.html] ... \end{browser}
%   \end{affianco}
% Le due colonne sono allineate in alto e formano un blocco unico, che non
% viene mai diviso fra due pagine. Con il secondo argomento facoltativo [c]
% le colonne sono centrate in verticale l'una rispetto all'altra (utile
% quando a destra c'è uno schema più basso del codice).
\newlength{\wtcodew}
\newlength{\wtresw}
\newlength{\wtgap}
\setlength{\wtgap}{0.5cm}
\NewDocumentEnvironment{affianco}{O{0.56} O{t}}{%
  \par\addvspace{6pt}\noindent
  \setlength{\wtcodew}{#1\linewidth}%
  \setlength{\wtresw}{\linewidth-\wtcodew-\wtgap}%
  % impostazioni valide in entrambe le colonne (vanno date prima della prima
  % minipage, altrimenti restano confinate in quella di sinistra)
  \lstset{xleftmargin=0pt, xrightmargin=0pt, aboveskip=0pt, belowskip=5pt}%
  \renewcommand{\wtbrowserbefore}{5pt}%
  \renewcommand{\wtbrowserafter}{6pt}%
  \renewcommand{\wtconsolebefore}{5pt}%
  \if c#2%
    \def\risultato{\end{minipage}\hspace{\wtgap}%
      \begin{minipage}[c]{\wtresw}}%
    \begin{minipage}[c]{\wtcodew}%
  \else
    \def\risultato{\end{minipage}\hspace{\wtgap}%
      \begin{minipage}[t]{\wtresw}\vspace{0pt}}%
    \begin{minipage}[t]{\wtcodew}\vspace{0pt}%
  \fi
}{%
  \end{minipage}\par\addvspace{10pt}}

% Listati un po' piu' piccoli, per le righe lunghe dentro una colonna stretta
\newcommand{\codicepiccolo}{%
  \lstset{basicstyle=\linespread{1.12}\scriptsize\FiraCodeNL}}

% Didascalia breve sotto un risultato (per esempio i dati inviati da un form)
\newcommand{\wtnota}[1]{%
  {\par\centering\Lato\footnotesize\color{neutral!30!black}#1\par}\vspace{4pt}}

% ============================================================
%  SCHEMA ACCANTO AL CODICE
% ============================================================
% Nella colonna di destra di "affianco":
%   \begin{immagine}[didascalia facoltativa]
%     \begin{tikzpicture} ... \end{tikzpicture}
%   \end{immagine}
% Lo schema viene centrato nella colonna e, se è più largo, ridotto alla sua
% larghezza (nel log compare "Schema ridotto", per accorgersene).
\newsavebox{\wtimgbox}
\NewDocumentEnvironment{immagine}{o}
  {\begin{lrbox}{\wtimgbox}}
  {\end{lrbox}%
   \par\noindent\makebox[\linewidth]{%
     \ifdim\wd\wtimgbox>\linewidth
       \typeout{Schema ridotto: \the\wd\wtimgbox\space -> \the\linewidth
                (riga \the\inputlineno)}%
       \resizebox{\linewidth}{!}{\usebox{\wtimgbox}}%
     \else
       \usebox{\wtimgbox}%
     \fi}\par
   \IfValueT{#1}{\vspace{4pt}\wtnota{#1}}}

% ============================================================
%  SCHEMI DELLA MEMORIA (variabili, oggetti, riferimenti, scope)
% ============================================================
% Livelli per gli scope annidati: il più esterno va disegnato sotto agli altri
\pgfdeclarelayer{liv1}
\pgfdeclarelayer{liv2}
\pgfdeclarelayer{liv3}
\pgfsetlayers{liv1,liv2,liv3,background,main}

\colorlet{mref}{accent!70!black}
\colorlet{mok}{green-gradient-6!85!black}
\colorlet{merr}{red-gradient-6!90!black}
\tikzset{
  % casella di una variabile: \node[mvar] (x) {42};
  mvar/.style={draw=neutral!55!black, fill=white, line width=0.7pt,
    minimum height=0.5cm, minimum width=1.4cm, inner sep=2pt,
    font=\FiraCodeNL\scriptsize},
  % nome della variabile, a sinistra della casella
  mnome/.style={font=\FiraCodeNL\scriptsize, text=primary!60!black,
    inner sep=2pt},
  % intestazione e righe di un oggetto (vedi \oggetto)
  mtesta/.style={draw=primary!65!black, fill=primary!16!white,
    line width=0.7pt, minimum height=0.46cm, inner sep=2pt,
    font=\Lato\scriptsize\bfseries, text=primary!35!black},
  mriga/.style={draw=primary!65!black, fill=primary!3!white,
    line width=0.7pt, minimum height=0.46cm, inner sep=3pt,
    font=\FiraCodeNL\scriptsize, align=left},
  % freccia di un riferimento e pallino da cui parte
  mrif/.style={-{Stealth[length=5pt,width=4pt]}, line width=0.85pt,
    draw=mref},
  mpunto/.style={circle, fill=mref, inner sep=1.4pt},
  % freccia "di lettura": dove il motore va a cercare un nome
  mcerca/.style={-{Stealth[length=5pt,width=4pt]}, line width=0.8pt,
    draw=green-gradient-6!80!black, dashed},
  % nota esplicativa nello schema
  mnota/.style={font=\Lato\scriptsize, text=neutral!20!black, align=left,
    inner sep=2pt},
  % codice dentro uno schema
  mcod/.style={font=\FiraCodeNL\scriptsize, inner sep=2pt},
  % riquadro di una funzione (oggetto funzione)
  mfun/.style={draw=accent!60!black, fill=accent!8!white, line width=0.7pt,
    rounded corners=6pt, minimum height=0.5cm, inner sep=3pt,
    font=\FiraCodeNL\scriptsize},
  % scope: \scopo{livello}{nome}{etichetta}{(nodo1) (nodo2) ...}
  mscopo1/.style={draw=neutral!50!white, fill=neutral!6!white,
    line width=0.7pt, rounded corners=5pt, inner sep=7pt},
  mscopo2/.style={draw=primary!45!white, fill=primary!6!white,
    line width=0.7pt, rounded corners=5pt, inner sep=7pt},
  mscopo3/.style={draw=accent!45!white, fill=accent!7!white,
    line width=0.7pt, rounded corners=5pt, inner sep=7pt},
  mscopolab/.style={font=\Lato\scriptsize\bfseries, text=neutral!25!black,
    inner sep=1.5pt},
}

% \oggetto[stile della testata]{nome}{(posizione)}{larghezza}{titolo}
%         {chiave/valore, chiave/valore, ...}
% Crea i nodi "nome" (testata), "nome-1", "nome-2", ... (righe) e "nome-box"
% (tutto l'oggetto). La posizione è quella del centro della testata, salvo
% un'ancora diversa nel primo argomento (per esempio anchor=north west).
% Senza righe ({}) disegna un oggetto vuoto. La larghezza della colonna delle
% chiavi si cambia con \mchiave{...} prima della chiamata.
% La larghezza indicata è quella minima: se una riga è più lunga, l'oggetto si
% allarga quanto serve (le righe non vanno mai a capo). Un valore che contiene
% una virgola va tra graffe: sep/{","}.
\newlength{\mchiavew}
\setlength{\mchiavew}{1.25cm}
\newcommand{\mchiave}[1]{\setlength{\mchiavew}{#1}}
\newlength{\mobjw}
\newlength{\mkeyw}
\newlength{\mrigaw}
\newsavebox{\wtimgmisura}
\NewDocumentCommand{\oggetto}{O{} m m m m m}{%
  \global\mobjw=#4\relax
  \global\mkeyw=\mchiavew
  \ifx\relax#6\relax\else
    % la colonna delle chiavi si allarga per la chiave più lunga
    \foreach \mk/\mv in {#6} {%
      \pgfinterruptpicture
        \global\setbox\wtimgmisura\hbox{\FiraCodeNL\scriptsize\mk}%
      \endpgfinterruptpicture
      \ifdim\dimexpr\wd\wtimgmisura+6pt\relax>\mkeyw
        \global\mkeyw=\dimexpr\wd\wtimgmisura+6pt\relax
      \fi}%
    \foreach \mk/\mv in {#6} {%
      % dentro tikzpicture il font è \nullfont: si misura fuori dalla figura
      \pgfinterruptpicture
        \global\setbox\wtimgmisura\hbox{\FiraCodeNL\scriptsize
          \makebox[\mkeyw][l]{\mk}\mv}%
      \endpgfinterruptpicture
      \setlength{\mrigaw}{\wd\wtimgmisura}%
      \ifdim\dimexpr\mrigaw+8pt\relax>\mobjw
        \global\mobjw=\dimexpr\mrigaw+8pt\relax
      \fi}%
  \fi
  \node[mtesta, minimum width=\mobjw, #1] (#2) at #3 {#5};
  \xdef\wtprev{#2}%
  \ifx\relax#6\relax
    \node[mriga, minimum width=\mobjw, anchor=north] (#2-1)
      at ([yshift=0.7pt]\wtprev.south) {\phantom{x}};
    \xdef\wtprev{#2-1}%
  \else
    \foreach \mk/\mv [count=\mi] in {#6} {%
      \node[mriga, minimum width=\mobjw, text width=\dimexpr\mobjw-6pt\relax,
            anchor=north] (#2-\mi) at ([yshift=0.7pt]\wtprev.south)
        {\makebox[\mkeyw][l]{\color{primary!55!black}\mk}\mv};
      \xdef\wtprev{#2-\mi}%
    }%
  \fi
  \node[fit=(#2)(\wtprev), inner sep=0pt] (#2-box) {};}

% \variabile[stile]{nome del nodo}{(posizione)}{nome mostrato}{valore}
% Il nome compare a sinistra della casella (nodo "nome-nome"). Per un
% riferimento si lascia il valore vuoto e si usa \riferimento.
\NewDocumentCommand{\variabile}{O{} m m m m}{%
  \node[mvar, #1] (#2) at #3 {#5};
  \node[mnome, anchor=east] (#2-nome) at ([xshift=-2pt]#2.west) {#4};}

% \riferimento[stile del percorso]{nodo di partenza}{percorso fino alla meta}
% es.: \riferimento{a}{-- (obj.west)}  oppure  \riferimento{a}{to[bend left] (obj)}
\NewDocumentCommand{\riferimento}{O{} m m}{%
  \draw[mrif, #1] (#2.center) #3;
  \node[mpunto] at (#2.center) {};}

% \scopo[stile]{livello 1-3}{nome}{etichetta}{(nodo) (nodo) ...}
% Il riquadro racchiude i nodi indicati e lascia in alto lo spazio per
% l'etichetta.
\NewDocumentCommand{\scopo}{O{} m m m m}{%
  \node[fit={#5}, inner sep=0pt] (#3-in) {};
  \node[mscopolab, anchor=south west] (#3-lab)
    at ([yshift=1pt]#3-in.north west) {#4};
  \begin{pgfonlayer}{liv#2}
    \node[mscopo#2, fit=(#3-in)(#3-lab), #1] (#3) {};
  \end{pgfonlayer}}

% \celle[stile]{nome}{(posizione)}{v0, v1, ...}: le caselle consecutive di un
% array, con l'indice sotto ciascuna. La posizione è il bordo sinistro della
% prima casella; crea i nodi nome-0, nome-1, ... e nome-box. Con \mvuota si
% disegna una casella vuota (un "buco"), con \noindici prima della chiamata si
% omettono gli indici.
\tikzset{
  mcella/.style={draw=primary!65!black, fill=primary!3!white, line width=0.7pt,
    minimum height=0.5cm, minimum width=0.75cm, inner xsep=4pt,
    font=\FiraCodeNL\scriptsize},
  mindice/.style={font=\FiraCodeNL\tiny, text=neutral!45!black, inner sep=1pt},
}
\newcommand{\mvuota}{{\Lato\itshape\color{neutral!50!black}vuota}}
\newif\ifmindici \mindicitrue
\newcommand{\noindici}{\mindicifalse}
\NewDocumentCommand{\celle}{O{} m m m}{%
  \coordinate (#2-inizio) at #3;
  \xdef\wtprev{#2-inizio}%
  \foreach \v [count=\i from 0] in {#4} {%
    \node[mcella, anchor=west, #1] (#2-\i)
      at ([xshift=-0.7pt]\wtprev.east |- #2-inizio) {\v};
    \ifmindici
      \node[mindice, anchor=north] at ([yshift=-1pt]#2-\i.south) {\i};
    \fi
    \xdef\wtprev{#2-\i}%
  }
  \node[fit=(#2-0)(\wtprev), inner sep=0pt] (#2-box) {};}

% Nodi di un albero DOM: elementi, testo, commenti
\tikzset{
  mdom/.style={rounded corners=2pt, draw=primary!60!black, line width=0.7pt,
    fill=primary!8!white, inner sep=2.5pt, font=\FiraCodeNL\scriptsize},
  mtesto/.style={mdom, draw=green-gradient-6, fill=green-gradient-1!45!white},
  mcommento/.style={mdom, draw=neutral!50!white, fill=neutral!8!white,
    text=neutral!40!black},
  mnuovo/.style={mdom, draw=accent!60!black, fill=accent!14!white},
  mramo/.style={draw=neutral!50!black, line width=0.55pt},
}

% \wtdialogo[valore del campo]{messaggio}{pulsanti}: una finestra di dialogo
% del browser (alert, confirm, prompt). Con l'argomento facoltativo compare
% anche il campo di testo di prompt.
\newcommand{\wtdialogohost}{localhost:3000}
\NewDocumentCommand{\wtdialogo}{o m m}{%
  \tikz\node[draw=neutral!45!white, fill=white, rounded corners=3pt,
             line width=0.7pt, inner sep=6pt, text width=4.6cm, align=left,
             font=\Lato\scriptsize]{%
    {\bfseries \wtdialogohost{} dice}\\[3pt]
    #2\IfValueT{#1}{\\[4pt]\wtcampo[4.3cm]{#1}}\\[5pt]
    \mbox{}\hfill #3};}

% \dbtabella[opzioni]{nome}{(posizione)}{nome della tabella}{n. colonne}{righe}
% Una tabella di database: la posizione è l'angolo in alto a sinistra, le righe
% si scrivono come in tabular (la prima è l'intestazione, seguita da \hline).
\NewDocumentCommand{\dbtabella}{O{} m m m m m}{%
  \node[inner sep=0pt, anchor=north west, #1] (#2) at #3 {%
    \FiraCodeNL\scriptsize\renewcommand{\arraystretch}{1.25}%
    \color{primary!55!black}%
    \begin{tabular}{|*{#5}{>{\color{black}}l|}}
      \hline #6 \\ \hline
    \end{tabular}};
  \node[anchor=south west, inner sep=1.5pt, font=\Lato\scriptsize\bfseries,
        text=primary!55!black] at (#2.north west) {#4};}

\usetikzlibrary{shadows}
% \wttoast[tipo]{titolo}{messaggio}: notifica della libreria ngx-toastr
% (tipo: success, error, warning, info), con i colori originali
\definecolor{toast-success}{HTML}{51A351}
\definecolor{toast-error}{HTML}{BD362F}
\definecolor{toast-warning}{HTML}{F89406}
\definecolor{toast-info}{HTML}{2F96B4}
\NewDocumentCommand{\wttoast}{O{info} m m}{%
  \tikz\node[fill=toast-#1, text=white, rounded corners=3pt, inner sep=6pt,
             text width=5.4cm, align=left, font=\Lato\scriptsize,
             execute at begin node={\hyphenpenalty=10000},
             drop shadow={opacity=0.22, shadow xshift=0.7pt, shadow yshift=-0.9pt}]
    {\IfValueT{#2}{{\bfseries #2}\\[1pt]}#3};}

% Segni di riuscita e di errore negli schemi
\newcommand{\mok}{{\color{mok}\itick}}
\newcommand{\merr}{{\color{merr}\icross}}

% ============================================================
%  CONTROLLI DEI FORM, COME LI DISEGNA UN BROWSER
% ============================================================
% I controlli dei form usano il carattere di sistema, non quello della pagina.
\newcommand{\wtuifont}{\Lato\scriptsize}
% \wtcampo[larghezza][sfondo]{testo}: campo di testo
\NewDocumentCommand{\wtcampo}{O{3cm} O{white} m}{%
  \fcolorbox{neutral!55!white}{#2}{\makebox[#1][l]{%
    \rule[-2.5pt]{0pt}{9.5pt}{\wtuifont #3}}}}
% \wtbottone[sfondo][testo][bordo]{etichetta}
\NewDocumentCommand{\wtbottone}{O{neutral!10!white} O{black} O{neutral!55!white} m}{%
  \fcolorbox{#3}{#1}{\rule[-2.5pt]{0pt}{9.5pt}{\wtuifont\color{#2}#4}}}
% \wtcasella{0|1}: checkbox vuota o spuntata
\newcommand{\wtcasella}[1]{\tikz[baseline=-2.8pt]{%
  \ifnum#1=1
    \fill[rounded corners=1pt, blue!70!black] (-3.2pt,-3.2pt) rectangle (3.2pt,3.2pt);
    \draw[white, line width=1.1pt] (-1.9pt,-0.2pt) -- (-0.5pt,-1.7pt) -- (2pt,1.8pt);
  \else
    \draw[rounded corners=1pt, neutral!60!black, line width=0.6pt]
      (-3.2pt,-3.2pt) rectangle (3.2pt,3.2pt);
  \fi}}
% \wtradio{0|1}: pulsante radio vuoto o selezionato
\newcommand{\wtradio}[1]{\tikz[baseline=-2.8pt]{%
  \ifnum#1=1
    \draw[blue!70!black, line width=0.9pt, fill=white] (0,0) circle (3.2pt);
    \fill[blue!70!black] (0,0) circle (1.7pt);
  \else
    \draw[neutral!60!black, line width=0.6pt] (0,0) circle (3.2pt);
  \fi}}
% \wttendina[larghezza]{valore}: menu a tendina chiuso
\NewDocumentCommand{\wttendina}{O{2.4cm} m}{%
  \fcolorbox{neutral!55!white}{neutral!6!white}{\makebox[#1][l]{%
    \rule[-2.5pt]{0pt}{9.5pt}{\wtuifont #2}\hfill{\tiny$\blacktriangledown$}}}}
% Puntatore del mouse, per indicare lo stato :hover
\newcommand{\wtpuntatore}{\tikz[baseline=0pt, scale=0.55]{%
  \filldraw[fill=white, draw=black, line width=0.35pt, line join=round]
    (0,0) -- (0,-0.62) -- (0.15,-0.48) -- (0.26,-0.72) -- (0.35,-0.68)
    -- (0.24,-0.44) -- (0.44,-0.44) -- cycle;}}
% Etichetta arrotondata a fianco di un elemento ("mouse sopra", "fuoco"...)
\newcommand{\wtstato}[1]{{\Lato\tiny\color{accent!75!black}\,(#1)}}

% ============================================================
%  PANNELLI DEI DEVTOOLS
% ============================================================
\definecolor{dtblue}{HTML}{0060DF}
\definecolor{dtbar}{HTML}{F0F0F4}
\definecolor{dtline}{HTML}{E0E0E6}
\definecolor{dtprop}{HTML}{0074E8}
\definecolor{dtsel}{HTML}{DD00A9}
\definecolor{dtok}{HTML}{2A8F3C}
\definecolor{dtwarn}{HTML}{B25E00}

% Barra delle schede dei DevTools; l'argomento e' la scheda attiva.
\newcommand{\dtscheda}[2]{%
  \ifthenelse{\equal{#1}{#2}}%
    {\tikz[baseline=(t.base)]{\node[inner sep=0pt, text=dtblue] (t) {\textbf{#2}};
      \draw[dtblue, line width=0.9pt] ([yshift=-2.5pt]t.south west)
        -- ([yshift=-2.5pt]t.south east);}}{#2}}
\newcommand{\dtschede}[1]{{\Lato\scriptsize\color{neutral!20!black}%
  \dtscheda{#1}{Analisi pagina}\hspace{9pt}\dtscheda{#1}{Console}\hspace{9pt}%
  \dtscheda{#1}{Debugger}\hspace{9pt}\dtscheda{#1}{Rete}\hspace{9pt}%
  \dtscheda{#1}{Editor stili}}}

% Finestra dei DevTools. Uso: \begin{devtools}{Rete} ... \end{devtools}
\newtcolorbox{devtools}[2][]{%
  enhanced, colback=white, colframe=neutral!45!white, boxrule=0.8pt, arc=3pt,
  left=0pt, right=0pt, top=0pt, bottom=0pt, boxsep=0pt,
  before skip=0.35cm, after skip=0.35cm,
  fonttitle=\scriptsize, colbacktitle=dtbar, coltitle=neutral!20!black,
  toptitle=3pt, bottomtitle=3pt, lefttitle=8pt,
  title={\dtschede{#2}}, #1}

% Pannello "Regole" (stili) dell'ispettore
\newcommand{\dtfont}{\FiraCodeNL\scriptsize}
% Risultato di console.table, dentro \begin{devtools}{Console}:
%   \ctabella{name & 'Hannibal' \\ age & 7}
\newcommand{\ctabella}[1]{%
  \par\vspace{6pt}\noindent\hspace*{8pt}%
  {\dtfont\renewcommand{\arraystretch}{1.3}\color{neutral!45!white}%
   \begin{tabular}{|>{\color{neutral!15!black}}l|>{\color{codestring}}l|}
     \hline
     \multicolumn{1}{|>{\color{neutral!15!black}}l|}{\textbf{(index)}} &
     \multicolumn{1}{>{\color{neutral!15!black}}l|}{\textbf{Values}} \\
     \hline
     #1 \\
     \hline
   \end{tabular}}\par\vspace{6pt}}
% \dtregola{selettore}{origine}: intestazione di una regola
\newcommand{\dtregola}[2]{%
  \par\noindent\hspace*{6pt}{\dtfont\textcolor{dtsel}{#1} \{}%
  \hfill{\Lato\tiny\color{neutral!45!black}#2}\hspace*{6pt}\par}
% \dtdich{proprieta}{valore} e \dtbarrata{proprieta}{valore}
\newcommand{\dtdich}[2]{%
  \par\noindent\hspace*{18pt}{\dtfont\textcolor{dtprop}{#1}: #2;}\par}
\newcommand{\dtbarrata}[2]{%
  \par\noindent\hspace*{18pt}{\dtfont\color{neutral!55!white}%
  \sout{#1: #2;}}\par}
\newcommand{\dtchiudi}{\par\noindent\hspace*{6pt}{\dtfont\}}\par\vspace{2pt}%
  {\color{dtline}\hrule}\vspace{3pt}}
% Intestazione di gruppo, per esempio "Ereditato da p"
\newcommand{\dtgruppo}[1]{%
  \par\vspace{1pt}\noindent\colorbox{dtbar}{\makebox[\dimexpr\linewidth-2\fboxsep][l]{%
    \Lato\tiny\color{neutral!25!black}#1}}\par\vspace{2pt}}

% ============================================================
%  COLORI CSS CON IL LORO VALORE ESATTO
% ============================================================
% Nei risultati "browser" si usano i colori con nome del CSS, che non
% coincidono con quelli omonimi di xcolor (ridefiniti dalla classe).
\definecolor{css-red}{HTML}{FF0000}
\definecolor{css-blue}{HTML}{0000FF}
\definecolor{css-green}{HTML}{008000}
\definecolor{css-yellow}{HTML}{FFFF00}
\definecolor{css-cyan}{HTML}{00FFFF}
\definecolor{css-hotpink}{HTML}{FF69B4}
\definecolor{css-teal}{HTML}{008080}
\definecolor{css-fuchsia}{HTML}{FF00FF}
\definecolor{css-darkred}{HTML}{8B0000}
\definecolor{css-deeppink}{HTML}{FF1493}
\definecolor{css-crimson}{HTML}{DC143C}
\definecolor{css-pink}{HTML}{FFC0CB}
\definecolor{css-darkblue}{HTML}{00008B}
\definecolor{css-aliceblue}{HTML}{F0F8FF}
\definecolor{css-link}{HTML}{0000EE}      % colore predefinito dei link
\definecolor{css-visited}{HTML}{551A8B}   % link visitati (Chrome/Firefox)

% Link come lo disegna il browser senza CSS: blu e sottolineato
\newcommand{\wtlink}[1]{\textcolor{css-link}{\uline{#1}}}
% \wtsopra{colore}{testo}: testo sopralineato (text-decoration: overline)
\newcommand{\wtsopra}[2]{\tikz[baseline=(t.base)]{%
  \node[inner sep=0pt, text=#1] (t) {\strut #2};
  \draw[#1, line width=0.45pt] ([yshift=-1.2pt]t.north west) -- ([yshift=-1.2pt]t.north east);}}

% ============================================================
%  FIELDSET (riquadro con legenda) come nei browser
% ============================================================
\newtcolorbox{wtfieldset}[1]{%
  enhanced, colback=white, colframe=neutral!45!white, boxrule=0.7pt,
  sharp corners, left=6pt, right=6pt, top=7pt, bottom=4pt,
  before skip=6pt, after skip=4pt, fontupper=\Lora\small,
  attach boxed title to top left={xshift=5pt, yshift=-2.3mm},
  boxed title style={colback=white, colframe=white, boxrule=0pt,
                     left=2pt, right=2pt, top=0pt, bottom=0pt},
  coltitle=black, fonttitle=\Lora\small, title={#1}}

% ============================================================
%  RICHIESTE NEL PANNELLO "RETE"
% ============================================================
% \dtstato{200}  \dtstato[dtwarn]{304}  : codice di stato colorato
\newcommand{\dtstato}[2][dtok]{%
  {\setlength{\fboxsep}{1.2pt}\colorbox{#1}{\FiraCodeNL\tiny\bfseries\color{white}#2}}}

% ============================================================
%  LE SPIEGAZIONI NON SI SPEZZANO FRA DUE PAGINE
% ============================================================
% I riquadri della classe (info, warning, error, definizione) sono dichiarati
% "breakable": quando non entravano nello spazio rimasto venivano tagliati a
% meta' frase, senza bordo inferiore, e proseguivano nella pagina dopo senza
% titolo. Qui la chiave "breakable" di tcolorbox viene ridefinita in modo che
% equivalga a "unbreakable": un riquadro che non entra passa intero alla pagina
% successiva, esattamente come i listati di codice.
\tcbset{breakable/.style={breakable@false}}


% ============================================================
%  Un concetto non si spezza fra due pagine
% ============================================================
% Il testo che introduce un listato, una figura, una lista o un risultato
% ("... come segue:") resta attaccato a ciò che introduce: se non entrano
% insieme in fondo alla pagina, passano insieme alla successiva. Le penalità
% impediscono il salto pagina solo in quel punto: il blocco che ne risulta
% (frase + listato, frase + figura, ...) è sempre più corto di una pagina.
\raggedbottom
% Un paragrafo non si divide mai fra due pagine
\interlinepenalty=10000
% Un codice inline lungo in fondo alla riga non esce dal margine: se il
% paragrafo non si spezza bene, TeX allarga un po' gli spazi di quella riga
\emergencystretch=3em
\makeatletter
% Niente salto pagina fra la frase introduttiva e la lista o il blocco
% centrato (tabelle, figure non flottanti) che la segue
\@beginparpenalty=10000
% Una lista resta tutta sulla stessa pagina
\@itempenalty=10000
\makeatother
\newcommand{\wtattacca}{\par\nopagebreak}
% (per i listati \wtattacca è già nell'hook che li chiude in una minipage)
\AddToHook{env/figure/before}{\wtattacca}
\AddToHook{env/table/before}{\wtattacca}
\AddToHook{env/console/before}{\wtattacca}
\AddToHook{env/browser/before}{\wtattacca}
\AddToHook{env/devtools/before}{\wtattacca}
\AddToHook{env/affianco/before}{\wtattacca}
% Un titolo non resta da solo in fondo alla pagina: se sotto non c'è spazio
% per qualche riga di contenuto, va a pagina nuova
%
% Inoltre ogni "blocco" (da un titolo di sezione o sottosezione al titolo
% successivo) che sta in una pagina non viene spezzato: la posizione di inizio
% e di fine di ogni blocco viene salvata nel file .aux (zref-savepos) e, alla
% compilazione successiva, prima del titolo si chiede tanto spazio quanto è alto
% il blocco. Se non c'è, il blocco intero passa alla pagina dopo. Servono quindi
% almeno tre compilazioni perché le posizioni si stabilizzino.
\usepackage{zref-savepos,zref-abspage}
\newcount\wtblocco
\makeatletter
\zref@addprop{savepos}{abspage}
\newcommand{\wtfineblocco}{%
  \ifnum\wtblocco>\z@ \zsavepos{wtb\the\wtblocco e}\fi}
\newcommand{\wtiniziablocco}{%
  \global\advance\wtblocco\@ne \zsavepos{wtb\the\wtblocco s}}
% \wtspazio{minimo}: Needspace pari all'altezza del prossimo blocco, se sta in
% una pagina, altrimenti pari al minimo. L'altezza misurata quando il blocco
% stava tutto in una pagina viene ricordata (\wtmem nel file .aux): così un
% blocco già spostato non torna indietro alla compilazione successiva.
\newdimen\wtaltezza
\newcommand{\wtmem}[2]{\expandafter\gdef\csname wtmem@#1\endcsname{#2}}
% \wtcap{n}: il blocco n apre un capitolo, quindi il blocco n-1 chiude il
% precedente e la sua "fine" viene registrata nella pagina del nuovo capitolo
\newcommand{\wtcap}[1]{\expandafter\gdef\csname wtcap@#1\endcsname{}}
% \wtsegue{n}: il blocco n contiene solo il titolo di una sezione, seguito
% subito da una sottosezione. Il titolo deve allora viaggiare insieme al blocco
% successivo: altrimenti, se la sottosezione passa alla pagina dopo, il titolo
% della sezione resta da solo in fondo alla pagina.
\newcommand{\wtsegue}[1]{\expandafter\gdef\csname wtsegue@#1\endcsname{}}
\newcommand{\wtspazio}[1]{%
  \begingroup
  \edef\wtn{\the\numexpr\wtblocco+1\relax}%
  \edef\wtps{\zref@extractdefault{wtb\wtn s}{abspage}{0}}%
  \edef\wtpe{\zref@extractdefault{wtb\wtn e}{abspage}{0}}%
  \wtaltezza=#1\relax
  \ifcsname wtsegue@\wtn\endcsname
    % solo un titolo: serve lo spazio del blocco che lo segue, piu' il titolo
    \edef\wtm{\the\numexpr\wtn+1\relax}%
    \edef\wtqs{\zref@extractdefault{wtb\wtm s}{abspage}{0}}%
    \edef\wtqe{\zref@extractdefault{wtb\wtm e}{abspage}{0}}%
    \ifnum\wtqs>\z@
      \ifnum\wtqe=\wtqs
        \wtaltezza=\dimexpr\zposy{wtb\wtm s}sp-\zposy{wtb\wtm e}sp
                   +4\baselineskip\relax
      \else\ifcsname wtmem@\wtm\endcsname
        \ifdim\csname wtmem@\wtm\endcsname<\maxdimen
          \wtaltezza=\dimexpr\csname wtmem@\wtm\endcsname+4\baselineskip\relax
        \fi
      \fi\fi
    \fi
  \else
  \ifnum\wtps>\z@
    \ifnum\wtpe=\wtps
      % tutto in una pagina: altezza esatta, da ricordare
      \wtaltezza=\dimexpr\zposy{wtb\wtn s}sp-\zposy{wtb\wtn e}sp\relax
      % si tiene la misura piu' alta: in cima alla pagina il blocco risulta piu'
      % basso (sparisce lo spazio prima del titolo), e senza questo l'impaginazione
      % oscillerebbe fra due soluzioni
      \ifcsname wtmem@\wtn\endcsname
        \ifdim\csname wtmem@\wtn\endcsname>\wtaltezza
          \wtaltezza=\csname wtmem@\wtn\endcsname\relax
        \fi
      \fi
      \immediate\write\@auxout{\string\wtmem{\wtn}{\the\wtaltezza}}%
    \else\ifnum\wtpe=\numexpr\wtps+1\relax
      \ifcsname wtcap@\the\numexpr\wtn+1\relax\endcsname
        % ultimo blocco di un capitolo: la fine e' nella pagina dopo
      \else\ifdim\zposy{wtb\wtn e}sp>\wttesta
        % la fine e' solo scivolata in cima alla pagina dopo: il blocco sta gia'
      \else\ifcsname wtmem@\wtn\endcsname
        \wtaltezza=\csname wtmem@\wtn\endcsname\relax
        \immediate\write\@auxout{\string\wtmem{\wtn}{\the\wtaltezza}}%
      \else\ifdim\zposy{wtb\wtn s}sp>\wttesta
        % comincia gia' in cima a una pagina ed e' spezzato: e' troppo lungo
        \immediate\write\@auxout{\string\wtmem{\wtn}{\the\maxdimen}}%
      \else
        % spezzato: si prova a portarlo intero alla pagina successiva
        \wtaltezza=\maxdimen
      \fi\fi\fi\fi
    \fi\fi
  \fi
  \fi
  \ifdim\wtaltezza=\maxdimen
    \par\newpage
  \else
    \ifdim\wtaltezza<#1\relax \wtaltezza=#1\relax\fi
    \ifdim\wtaltezza>.97\textheight \wtaltezza=#1\relax\fi
    \Needspace{\dimexpr\wtaltezza+\baselineskip\relax}%
  \fi
  \endgroup}
% posizione verticale oltre la quale un punto e' "in cima alla pagina"
\newdimen\wttesta
\AtBeginDocument{\wttesta=\dimexpr\paperheight-\topmargin-1in-\headheight-\headsep-40pt\relax}
\AddToHook{cmd/chapter/before}{\wtfineblocco\wtiniziablocco\immediate\write\@auxout{\string\wtcap{\the\wtblocco}}}
\AddToHook{cmd/section/before}{\wtfineblocco\wtspazio{22\baselineskip}\wtiniziablocco}
\AddToHook{cmd/subsection/before}{\wtfineblocco
  \if@nobreak\immediate\write\@auxout{\string\wtsegue{\the\wtblocco}}\fi
  \wtspazio{8\baselineskip}\wtiniziablocco}
\AddToHook{enddocument}{\wtfineblocco}
\makeatother
